% Compile with XeLaTeX (or pdfLaTeX with the ko.TeX/kotex package installed).
% Korean support via kotex.
\documentclass[10pt,conference]{IEEEtran}

\usepackage{kotex}            % Korean
\usepackage{cite}
\usepackage{amsmath,amssymb,amsfonts}
\usepackage{algorithmic}
\usepackage{graphicx}
\usepackage{booktabs}
\usepackage{textcomp}
\usepackage{xcolor}
\usepackage{listings}
\usepackage{url}

\lstset{basicstyle=\ttfamily\footnotesize,breaklines=true,frame=single,
        columns=fullflexible,keepspaces=true}

\begin{document}

\title{ChatGPT가 만든 버그는 다르다:\\
데이터 누수 없는 환경에서 LLM 생성 코드 결함의\\
분류 체계와 체계성(Systematicity) 분석}

\author{\IEEEauthorblockN{(저자명)}
\IEEEauthorblockA{성균관대학교 소프트웨어학과\\
이메일: dwchoi95@g.skku.edu}}

\maketitle

\begin{abstract}
학생과 개발자가 대규모 언어 모델(LLM)이 생성한 코드를 제출하는 사례가 급증하면서,
\emph{LLM이 만든 버그는 인간이 만든 버그와 다른가, 다르다면 기존 자동 프로그램
수리(APR) 도구로 효과적으로 고쳐지는가}라는 질문이 중요해지고 있다. 그러나 이 질문을
실증적으로 답하는 데에는 방법론적 딜레마가 있다. 모델의 지식 컷오프 \emph{이후} 문제를
쓰면 데이터 누수(leakage)는 피하지만 그 시기의 인간 제출은 그 자체가 AI 코드일 수 있어
인간 대조군이 오염되고, \emph{이전} 문제를 쓰면 인간 코드는 진정하지만 누수가 발생한다.
본 연구는 이 두 제약이 \emph{서로 다른 산출물}에 \emph{서로 다른 컷오프}로 걸린다는 점에
착안해, 모델 컷오프(\texttt{gpt-3.5-turbo}, 2021년 9월)와 ChatGPT 출시(2022년 11월)
\emph{사이의 겹치는 구간}을 활용한다. 이 구간의 AtCoder 문제는 모델이 학습하지 못했고
(누수 없음), 그 시기 인간 제출은 ChatGPT 이전이라 진정하다. ConDefects를 기반으로
428개 문제에 대해 4{,}280개의 솔루션을 생성하고 전체 숨은 테스트로 채점한 뒤, 개방
코딩(open coding)으로 데이터-주도 결함 분류 체계를 도출하였다. 분석 결과, LLM 결함의
76\%가 \textbf{명세 오해(Spec-Misinterpretation)}로, \emph{문법적으로 정상이고 실행되지만
미묘하게 다른 문제를 푸는} 코드였으며, API 환각 같은 고전적 범주는 거의 나타나지 않았다.
나아가 한 문제에 대한 여러 생성물의 버그가 같은 실패 유형으로 강하게 수렴함을
정량적으로 보였다(평균 최빈비율 0.845, 문제의 38.8\%가 완전 일치). 이는 LLM 버그가
무작위가 아니라 \emph{체계적}이며, 따라서 명세를 두 번째 오라클로 삼는
명세-기반 수리(Spec-Grounded Repair)가 유망함을 시사한다. 누수 없고 인간과 짝지어진
데이터셋과 방법론을 공개한다.
\end{abstract}

\begin{IEEEkeywords}
LLM 생성 코드, 버그 분류 체계, 자동 프로그램 수리, 데이터 누수, 명세-기반 수리
\end{IEEEkeywords}

\section{서론}
ChatGPT를 비롯한 대규모 언어 모델(LLM)이 코드를 생성하는 능력은 교육과 산업 현장 모두에서
빠르게 확산되었다~\cite{chen2021codex}. 특히 프로그래밍 교육에서는 학생이 과제 코드를
LLM에게 생성시켜 제출하는 일이 흔해졌고, 이로 인해 \emph{LLM이 만든 코드 고유의 버그
패턴}이라는 새로운 연구 대상이 등장하였다. 기존 자동 프로그램 수리(APR) 도구는 대체로
\emph{인간이 작성한} 버그를 전제로 설계되었기 때문에~\cite{monperrus2018survey,hu2019refactory},
LLM 버그가 인간 버그와 다른 분포를 가진다면 기존 도구의 효과가 떨어질 수 있다.

이 가설을 실증적으로 검증하려면 \emph{진정한} 데이터가 필요하지만, 두 가지 요구가 충돌한다.
첫째, AI 생성 코드가 \emph{진짜 생성}이려면 대상 문제가 모델의 지식 컷오프 \emph{이후}에
출제되어 모델이 정답을 암기하지 않아야 한다(데이터 누수 회피~\cite{wu2023condefects}).
둘째, 비교 대상인 \emph{인간} 코드가 진정하려면 LLM 대중화 \emph{이전}에 작성되어야 한다.
컷오프 이후 시기에는 인간 제출조차 AI 보조일 수 있어 대조군이 오염되기 때문이다. 전자는
``이후''를, 후자는 ``이전''을 요구하므로, 같은 문제 집합에서 두 조건을 동시에 만족하기
어렵다.

본 연구의 핵심 통찰은 \textbf{두 제약이 서로 다른 산출물에 서로 다른 컷오프로 걸린다}는
것이다. 누수는 \emph{생성 대상 문제}에 걸리고(모델 컷오프 기준), 인간 진정성은
\emph{인간 제출}에 걸린다(LLM 대중화 기준). 두 컷오프가 다르므로, 그 사이에 \emph{두
조건이 동시에 성립하는 겹치는 구간}이 존재한다. 우리는 \texttt{gpt-3.5-turbo}의 지식 컷오프
(2021년 9월)와 ChatGPT 출시(2022년 11월 30일) 사이의 구간을 활용한다. 이 구간에 출제된
AtCoder 문제는 모델이 학습하지 못했고(누수 없음), 그 컨테스트의 인간 제출은 ChatGPT
이전이므로 진정하다.

본 논문의 기여는 다음과 같다.
\begin{itemize}
\item \textbf{누수 없고 인간-진정한 실험 설계.} 컷오프--출시 겹치는 구간을 이용해 위 딜레마를
해소하고, 같은 문제에서 AI와 인간 결함을 짝지어 비교할 수 있게 한다.
\item \textbf{데이터-주도 결함 분류 체계.} 428개 문제·4{,}280개 생성물의 실패를 개방 코딩으로
분류하여, 명세 오해(Spec-Misinterpretation)가 76\%로 지배적이며 API 환각은 거의 없음을
보인다.
\item \textbf{체계성(systematicity)의 정량화.} 한 문제의 여러 생성 버그가 같은 실패 유형으로
수렴함을 보여(평균 최빈비율 0.845), LLM 버그가 무작위가 아닌 \emph{체계적} 현상임을
입증한다.
\item \textbf{공개 데이터셋과 방법론.} 문제 명세, 전체 숨은 테스트, AI 버그(문제당 10개),
짝지어진 인간 버그(결함 위치 포함)를 공개한다.
\end{itemize}

\section{배경 및 관련 연구}
\subsection{LLM 코드 생성과 데이터 누수}
Codex와 그 평가 벤치마크(HumanEval)~\cite{chen2021codex} 이후, LLM 기반 코드 생성과 수리는
빠르게 발전하였다~\cite{xia2023llmapr}. 그러나 공개 벤치마크가 학습 데이터에 포함되는
\emph{데이터 누수} 문제는 평가의 타당성을 훼손한다. ConDefects~\cite{wu2023condefects}는
모델 컷오프 이후(2021년 10월\,--\,)에 출제된 AtCoder 문제만을 모아 이 문제를 완화한
결함 위치추적·수리용 데이터셋이다. 본 연구는 ConDefects를 토대로 하되, \emph{인간 대조군
진정성}이라는 추가 제약을 시기 윈도우로 함께 다룬다.

\subsection{자동 프로그램 수리(APR)와 교육}
전통적 APR은 탐색 기반·의미 기반·학습 기반으로 발전해 왔으며~\cite{monperrus2018survey},
최근에는 LLM을 직접 수리에 활용한다~\cite{xia2023llmapr}. 교육용 입문 프로그래밍에서는
Refactory~\cite{hu2019refactory}가 정답 풀이 군을 재구성하여 학생 버그를 수리한다. 이들
도구는 모두 인간 버그를 전제로 한다는 공통점이 있다.

\subsection{경쟁 프로그래밍 데이터}
IBM Project CodeNet~\cite{puri2021codenet}과 Codeforces 계열 데이터는 인간 제출을 풍부히
제공하나, 모두 LLM 대중화 \emph{이전}에 구축되어 누수 문제에서 자유롭지 못하다. 본 연구는
의도적으로 누수가 없는 ConDefects(AtCoder)를 택한다.

\section{연구 방법}
\subsection{시기 윈도우 설계}
생성 모델로 \texttt{gpt-3.5-turbo}(지식 컷오프 2021년 9월)를 고정한다. 데이터 누수가 없으려면
대상 문제가 2021년 9월 이후 출제여야 하고, 인간 대조군이 진정하려면 ChatGPT 출시
(2022년 11월 30일) \emph{이전} 컨테스트여야 한다. 따라서 \textbf{[2021-09, 2022-11]} 구간의
문제만 사용한다. 이 구간에서 ConDefects의 Python 과제 중 인간 정답·버그 짝과 전체 숨은
테스트를 갖춘 428개 문제를 얻는다(인터랙티브 문제 4개 제외). 짝지어진 인간 버그(결함 위치
포함)는 1{,}215개이다.

\subsection{명세, 테스트, 생성}
ConDefects는 코드와 테스트는 제공하나 자연어 명세는 포함하지 않으므로, 각 과제의 영어
명세를 AtCoder에서 수집해 구조적으로 정규화하였다. 채점에는 ConDefects의 \emph{전체 숨은
테스트}(문제당 평균 약 32개)를 사용한다. 생성은 다음 설정으로 수행한다: 모델
\texttt{gpt-3.5-turbo}, 온도 1.0, 문제당 10개. 프롬프트는 ``Python 3 프로그램만 출력하라''는
간결한 시스템 지시와 정규화된 명세(사용자 메시지)로 구성한다. 이 시스템 지시 한 줄이
없으면 생성물의 46\%가 C++ 등 비-Python으로 이탈하므로, 본 설정에서 모든 생성물은
100\% Python이다.

\subsection{분류 방법}
사전 범주를 두지 않고 실패 코드를 개방 코딩으로 분류한다. 특히 오답(Wrong Answer; WA)은
실행은 되나 출력이 틀린 경우로, \emph{무엇을} 오해했는지 알기 위해 각 AI 버그를 같은 문제의
인간 정답(\texttt{correct.py})과 대조하였다. 신뢰도는 카테고리 확정 후 두 저자의 독립 분류로
Cohen's $\kappa$~\cite{cohen1960kappa}를 측정한다(목표 $\kappa\geq 0.7$).

\subsection{체계성 지표}
LLM 버그가 무작위인지 체계적인지 평가하기 위해, 문제별로 그 버그들의 실패 유형 분포를
본다. 문제 $p$의 버그 집합에서 가장 흔한 유형의 비율을 \emph{최빈비율(modal-share)},
유형 분포의 정규화 섀넌 엔트로피를 \emph{비체계성}으로 정의한다. 최빈비율이 1이면 모든
버그가 같은 유형(완전 체계적)이다.

\section{결과}
\subsection{채점 결과}
4{,}280개 생성물 중 589개(13.8\%)만이 전체 숨은 테스트를 통과하였고, 3{,}691개가 버그였다.
표~\ref{tab:difficulty}에서 보듯 통과율은 난이도에 따라 급격히 떨어져, 어려운 문제에서는
사실상 0\%였다. 주목할 점은 공개 샘플 입출력만으로 채점했을 때의 통과율이 23\%로
\emph{과대평가}되었다는 것이다. 샘플만 통과하고 숨은 테스트에서 떨어지는 약 200개 솔루션이
보정되었으며, 이는 LLM 코드 평가에 전체 테스트가 필수임을 보여준다.

\begin{table}[t]
\caption{난이도별 통과율 (전체 숨은 테스트, 문제당 10생성)}
\label{tab:difficulty}
\centering
\begin{tabular}{lrrr}
\toprule
난이도(AtCoder) & 생성 & 통과 & 통과율 \\
\midrule
쉬움 ($<$400)        & 1{,}200 & 550 & 45.8\% \\
중간 (400--1200)     & 1{,}040 & 39  & 3.7\% \\
어려움 (1200--2000)  & 890     & 0   & 0.0\% \\
매우 어려움 ($\geq$2000) & 1{,}150 & 0 & 0.0\% \\
\midrule
합계                 & 4{,}280 & 589 & 13.8\% \\
\bottomrule
\end{tabular}
\end{table}

\subsection{데이터-주도 결함 분류}
표~\ref{tab:taxonomy}는 도출된 분류 분포이다. \textbf{명세 오해(WA)}가 전체 버그의 76\%로
압도적이며, 비효율로 인한 시간 초과(TLE) 14\%, 런타임 오류(RE) 9\%, 출력 잘림 등(CE)
1\%가 뒤따른다. 특히 \emph{API 환각(Hallucinated-API)은 사실상 나타나지 않았다}. 관찰된
\texttt{ModuleNotFoundError}는 존재하지 않는 함수 호출이 아니라 \texttt{sympy} 등 채점
샌드박스에 미설치된 \emph{실제} 라이브러리 때문이었다. 이는 선행 제안의 5범주가 모든
설정에 보편적이지 않으며 데이터-주도 분류가 정당함을 시사한다.

\begin{table}[t]
\caption{AI 버그 분류 분포 (총 3{,}691 버그)}
\label{tab:taxonomy}
\centering
\begin{tabular}{lrr}
\toprule
카테고리 & 수 & 비율 \\
\midrule
명세 오해 / 불완전 로직 (WA) & 2{,}808 & 76\% \\
비효율 / 복잡도 (TLE)       & 504 & 14\% \\
런타임 오류 (인덱스/입력형식/수치) & 338 & 9\% \\
출력 잘림 / 형식 오류 (CE)  & 41 & 1\% \\
\bottomrule
\end{tabular}
\end{table}

\subsection{명세 오해의 하위 유형}
인간 정답과의 대조를 통해 명세 오해를 다섯 하위 유형으로 세분하였다.
\begin{itemize}
\item \textbf{S1 출력/케이스 단순화:} 출력 분기를 하나의 규칙으로 뭉갬(예: ``er''/``ist''를
끝 두 글자로 출력).
\item \textbf{S2 핵심 조건 단순화:} 진짜 요구를 더 쉬운 대용 조건으로 대체(예: 흑칸 연결성을
``각 행에 칸이 있는가''로).
\item \textbf{S3 잘못된 알고리즘/공식:} 틀린 수학 모델을 적용(예: 중복 순열 수를 일반 순열
공식으로).
\item \textbf{S4 패러다임 미인식:} 어려운 문제에서 필요한 알고리즘을 인식하지 못하고 완전
탐색·임기응변으로 대체(예: 위상 정렬을 \texttt{permutations}로).
\item \textbf{S5 엣지/언어 의미:} 구조는 맞으나 경계에서 실패(예: \texttt{round()}의 은행가
반올림).
\end{itemize}
공통적으로 AI 버그는 장황하고 주석이 많아 \emph{그럴듯해 보이지만 틀린} 형태였으며, 이는
문법·패턴 기반 도구가 이를 놓치기 쉬운 이유와 부합한다.

\subsection{체계성}
표~\ref{tab:systematicity}는 체계성 지표이다. 버그가 2개 이상인 387개 문제에서 평균
최빈비율은 0.845로, 한 문제의 버그 중 평균 84.5\%가 같은 실패 유형이었다. 문제의
38.8\%(150/387)는 모든 버그가 단일 유형으로 완전히 일치하였다. 예컨대 반올림 문제에서는
대부분의 시도가 동일하게 \texttt{round()}를 사용했고, 연결성 문제에서는 동일하게 칸 존재
여부로 단순화하였다. 즉 \textbf{LLM은 한 문제를 일관되게 같은 방식으로 오해}한다. 이 측정은
거친(coarse) 유형 기준의 \emph{하한}이며, 명세 오해의 하위 유형 기준으로는 더 강한 수렴이
관찰된다.

\begin{table}[t]
\caption{문제별 체계성 (버그 $\geq 2$인 387문제)}
\label{tab:systematicity}
\centering
\begin{tabular}{lr}
\toprule
지표 & 값 \\
\midrule
평균 최빈비율(modal-share)      & 0.845 \\
완전 일치 문제 비율              & 38.8\% \\
평균 정규화 엔트로피(0=완전체계) & 0.43 \\
최빈비율 = 1.0 인 문제 수        & 150 / 387 \\
\bottomrule
\end{tabular}
\end{table}

\section{논의}
\subsection{명세-기반 수리의 함의}
명세 오해가 지배적이고(76\%) 체계적(최빈비율 0.845)이라는 두 결과는 직접적인 수리 전략을
시사한다. 버그가 문법적으로 정상이고 실행되며 \emph{명세}만 어긋난다면, 자연어 명세에서
입력 범위·출력 형식·엣지 조건 등 의미적 제약을 추출하여 ``테스트 통과''와 별개의
\emph{두 번째 오라클}로 삼는 \textbf{명세-기반 수리}가 유효할 수 있다. 더욱이 한 문제의
오답들이 같은 오해로 수렴하므로, 그 오해를 한 번 교정하면 다수 오답을 일괄 수리할 여지가
있다.

\subsection{프롬프트와 언어 이탈}
언어를 명시하지 않으면 생성물의 46\%가 C++ 등으로 이탈하였다. 이는 평가 설계 시 생성
프롬프트가 결과 분포를 크게 좌우함을 보여주며, 본 연구는 ``Python 3 프로그램만 출력''
지시로 이를 통제하였다.

\subsection{타당성 위협}
\textbf{내적 타당성.} 단일 모델(\texttt{gpt-3.5-turbo})만을 사용하였으므로 다른 모델로의
일반화는 추후 과제이다. AtCoder의 일부 special-judge(복수 정답) 문제는 정확 일치 비교로
오판될 수 있으나 그 비중은 작다. \textbf{외적 타당성.} 인간 진정성은 ``컨테스트 시기가
ChatGPT 이전''이라는 \emph{프록시}에 의존한다(개별 제출 시각은 데이터셋에 없다). 옛 문제에
대한 후대 AI 재제출 가능성은 소수로 남는다. \textbf{구성 타당성.} 공개 샘플만으로 채점하면
통과율이 과대평가되므로(23\% vs 13.8\%) 전체 숨은 테스트를 사용하였다.

\section{결론}
본 연구는 모델 컷오프와 ChatGPT 출시 사이의 겹치는 구간을 활용하여, 데이터 누수가 없고
인간 대조군이 진정한 환경에서 LLM 생성 코드의 결함을 분석하였다. 결함의 76\%가 명세
오해였고 API 환각은 거의 없었으며, 한 문제에 대한 버그들이 같은 실패 유형으로 강하게
수렴하는 \emph{체계성}을 정량적으로 확인하였다. 이는 명세를 두 번째 오라클로 삼는
명세-기반 수리의 타당성을 뒷받침한다. 향후 연구로는 (1) 짝지어진 인간 버그와의 분포 비교
및 결함 위치 정확도 분석, (2) 기존 APR 도구의 범주별 수리율 측정, (3) 명세-기반 수리
프로토타입의 구현과 평가를 계획한다.

\section*{데이터 가용성}
명세·전체 테스트·AI 버그(문제당 10개)·짝지어진 인간 버그를 포함한 데이터셋과 코드는
\url{https://github.com/dwchoi95/aitaxo} 에서 공개한다.

\bibliographystyle{IEEEtran}
\bibliography{references}

\end{document}
